%!TEX root = ../../../exa-ma-deliverables.tex
\ApplicationFicheHeading{app-feelpp-thermomeca-1}
\label{sec:app:specs:app-feelpp-thermomeca-1}
\begin{FicheMetadataTable}
id & \nolinkurl{app-feelpp-thermomeca-1} \\
name & Thermomechanical coupling \\
Responsible & V. Chabannes; C. Prud'homme; I. Rami\`ere; H. Barucq \\
work\_package & WP1, WP3 \\
application\_type & demonstrator \\
purpose & Verify partitioned thermoelastic coupling using a displacement-dependent effective conductivity on a fixed mesh; assess solver performance and HPC scalability. \\
Framework & Feel++, MPI \\
operational status & implementation-in-progress \\
scientific reference & Dubois et al., CSMA 2026; \href{https://cea.hal.science/cea-05590574}{cea-05590574} \\
Feel++ configuration & Heat and solid models and coupling settings planned \\
benchmark data & Published Cast3M/MFEM study; Feel++ run archive pending \\
\end{FicheMetadataTable}

\subsection{Description of the benchmark}
This application studies static linear thermoelasticity with two-way coupling.
Thermal expansion changes the geometry seen by the heat equation, while temperature
drives the mechanical loading. The reference study by
\citeA{dubois_couplage_2026} formulates heat conduction on the initial geometry
using a displacement-dependent effective conductivity, verifies equivalence
with explicit mesh deformation in Cast3M, and studies parallel performance in
MFEM. It acknowledges Exa-MA support under ANR-22-EXNU-0002.
The article and its benchmark descriptions are publicly available through
\href{https://cea.hal.science/cea-05590574}{HAL (cea-05590574)} and the
\href{https://csma2026.sciencesconf.org/688246/document}{CSMA 2026 proceedings}.

\paragraph{Application scope and maturity.}
The Feel++ demonstrator adopts this study as a scientific reference for WP1
multiphysics coupling and WP3 solver/preconditioner evaluation. The published
Cast3M/MFEM measurements below are reference results, not completed Feel++ runs.
The dedicated Feel++ thermomechanics case is still in development under
\nolinkurl{BM-FEEL-THERMOMECA-01}.

\paragraph{Coupled model.}
Let $\Omega$ be the reference domain, $x=X+u(X)$ the deformation,
$F=I+\nabla u$ and $J=\det F$. In the small-strain, static setting of
\citeA{dubois_couplage_2026}, the reference-configuration problem is
\begin{gather}
 -\nabla\cdot(K(u)\nabla T)=J s,
 \qquad K(u)=JF^{-1}kF^{-T},\label{eq:app:thermomeca:heat}\\
 -\nabla\cdot\sigma(u,T)=B,
 \qquad\varepsilon(u)=\tfrac12(\nabla u+\nabla u^T),\label{eq:app:thermomeca:mechanics}\\
 \sigma(u,T)=\lambda\operatorname{tr}(\varepsilon(u))I
 +2\mu\varepsilon(u)-(3\lambda+2\mu)\alpha(T-T_{\mathrm{ref}})I.
 \label{eq:app:thermomeca:stress}
\end{gather}
Here $k$ and $s$ are conductivity and volumetric heat source on the deformed
domain; $K$ and $Js$ are their effective reference-domain counterparts.
Temperature is in K, displacement in m, conductivity in
$\mathrm{W\,m^{-1}\,K^{-1}}$, stress in Pa, heat source in
$\mathrm{W\,m^{-3}}$, and body force $B$ in $\mathrm{N\,m^{-3}}$.
The Piola heat flux $Q=JF^{-1}q$, with $q=-k\nabla_xT$, preserves integrated
normal heat flux. Neumann data must be transformed consistently using
$n\,da=JF^{-T}N\,dA$; homogeneous Neumann data remain homogeneous.
Dirichlet data are pulled back through $x=X+u(X)$.
The comparison formulation solves heat conduction on the explicitly deformed
mesh and mechanics on the reference mesh.

\paragraph{Reference cases.}
The bar follows \citeA{novascone_coupling_2015}; the second case is a fuel-pellet
fragment inspired by the ALCYONE setting \cite{introiini_alcyone_2024}.
\begin{table}[!ht]
\centering\small
\begin{tabularx}{\textwidth}{@{}>{\raggedright\arraybackslash}p{0.19\textwidth}>{\raggedright\arraybackslash}X>{\raggedright\arraybackslash}X@{}}
\toprule
\textbf{Parameter} & \textbf{Heated bar} & \textbf{Fuel-pellet fragment} \\
\midrule
Geometry & Length $5$ m; section $1\times1$ m$^2$ & Height $7$ mm; radius $4$ mm \\
Thermal loading & End temperatures $100$ K and $1100$ K; no volumetric source & Surface facing the cladding at $600$ K; volumetric heat source \\
Mechanical constraints & Axial deformation only & Constraints representing stacking and confinement; exact boundary masks to be supplied \\
Material data & Taken from the Novascone reference; exact input deck to be confirmed & $k=2.5I$ W\,m$^{-1}$\,K$^{-1}$; $E=190$ GPa; $\nu=0.3$; $\alpha=10^{-5}$ K$^{-1}$; $T_{\mathrm{ref}}=96.8$ K \\
Published meshes & $2560$ nodes for field plots; $100\,000$ for sequential timings; $83\times10^6$ for MFEM scaling & $120\,221$ nodes for Cast3M verification \\
\bottomrule
\end{tabularx}
\caption{Case data from Dubois et al. (2026), Sections 3.1--3.3. Node counts are not degree-of-freedom counts.}
\label{tab:app:thermomeca:cases}
\end{table}

\subsection{Benchmarking methodology}
\paragraph{Numerical verification (WP1).}
Use partitioned Gauss--Seidel iterations, solving the thermal and mechanical
problems successively. The paper stops when the relative change in the
Euclidean displacement norm falls below $10^{-6}$. Its exact residual
normalization must be recorded with the Feel++ input deck.
Compare both coupling formulations at identical mesh resolution, material data,
boundary conditions and solver tolerances. Compare temperature and displacement
after mapping to the same reference geometry, and report coupling iteration
counts and relative field errors. The paper reports near-roundoff agreement;
a Feel++ acceptance tolerance must be fixed before its runs.

\paragraph{Solver and scaling protocol (WP3).}
The MFEM reference compares MUMPS with conjugate gradients preconditioned by
HypreBoomerAMG (HMIS coarsening, strong threshold $0.25$, absolute residual
tolerance $10^{-6}$). It reuses the mechanical factorization for the direct
solver and the previous coupling iterate as the initial guess for the iterative
solver. These are reference settings, not verified Feel++ command-line options.
Measure thermal preparation/assembly, thermal solve, mechanical
preparation/assembly, mechanical solve and total coupling time separately.
Record linear and coupling iteration counts alongside timing.

For strong scaling, hold the global problem and accuracy fixed. The published
MFEM campaign uses $83\times10^6$ mesh nodes and $64$--$4096$ MPI processes,
one CPU core per process, on CCRT TOPAZE: AMD EPYC Milan 7763 at $2.45$ GHz,
$128$ cores and $256$ GB RAM per node. For a future Feel++ campaign,
report $S(p)=t(p_0)/t(p)$ and
$E(p)=S(p)/(p/p_0)$ with an explicitly stated baseline $p_0$, together with
repeated-run variability and all machine/software settings.

\subsection{Input/Output Dataset Description}
\paragraph{Available reference documentation.}
The open CSMA paper provides the formulation, case descriptions, sequential
timing Tables 1--2 and the TOPAZE strong-scaling plot (Figure 4).
The original bar study has a
\href{https://doi.org/10.1016/j.nucengdes.2015.07.005}{publisher record} and an
\href{https://www.osti.gov/servlets/purl/1248174}{open INL manuscript}.
The article documents the problem and published timings but does not supply
an executable benchmark package with the input decks, mesh generator and
raw scaling measurements.

\paragraph{Feel++ inputs and outputs to archive.}
The catalogue specifies Gmsh meshes and JSON configuration as inputs, with HDF5
fields and CSV outputs. Supply both reference cases with material values,
boundary labels, source amplitude, finite-element order, quadrature and solver
settings. Export $T$, $u$, coupling/linear residual histories, field-comparison
errors and phase timings. Record the software versions, launch commands,
mesh identifiers, dependency versions, MPI placement and repeated-run
measurements before validating the benchmark.

\subsection{Results summary}
\paragraph{Published reference results: Cast3M and MFEM.}
\citeA{dubois_couplage_2026} reports three coupling iterations for the bar and
five for the pellet. The Cast3M formulations agree to approximately $10^{-15}$
for the bar; maximum normalized discrepancies for the pellet are of order
$10^{-13}$ in temperature and $10^{-14}$ in displacement. The bar comparison
also agrees with the MOOSE/BISON reference solution.

On the $100\,000$-node sequential bar case, the reported total Cast3M times are
$249.188$ s for effective conductivity and $253.180$ s for explicit deformation.
MFEM totals are $62.26$ s with MUMPS and $72.23$ s with CG/AMG. Both timing
tables identify an Intel Core Ultra 5 135H machine (18 threads, 16 GB RAM);
these measurements are not a general cross-framework speedup claim.

For MFEM strong scaling, the mechanical solve accounts for over $90\%$ of total
time and scales nearly ideally through $1024$ MPI processes, with further gains
but reduced efficiency beyond that. The thermal solve scales well through
$2048$ processes, then its time rises slightly at $4096$; at $2048$ it contributes
less than $1\%$ of total time. Exact scaling timings are plotted in the paper,
not tabulated as an accompanying raw dataset.

\paragraph{Feel++ validation and benchmark results still required.}
No measured Feel++ verification error, timing or scalability result is claimed
for \nolinkurl{BM-FEEL-THERMOMECA-01}.
\paragraph{Planned contribution.} Obtain the authors' reference
input decks and raw data; confirm all bar material parameters, pellet source
and boundary masks, discrete spaces and stopping-test definition; supply reviewed heat/solid CFG and JSON models plus coupling settings, using a custom driver only if needed; select acceptance tolerances. A Feel++ run archive remains planned.

% Generated by scripts/feelpp_catalogue.py; edit data/feelpp-mesh-catalogue.json.

\subsection{Application package and shared meshes}
\begin{FicheMetadataTable}
Repository entry & \nolinkurl{apps/app-feelpp-thermomeca-1/} (local review layout in apps-feelpp; publication pending) \\
Delivery model & toolbox config planned \\
Available material & Detailed CSMA reference, equations and clearly attributed Cast3M/MFEM results. \\
Results status & Feel++ campaign results planned. \\
\hyperref[mesh:feelpp:thermoelastic-bar]{Heated thermoelastic bar} & 5 m x 1 m x 1 m bar. Supply geometry generator, end/constraint markers, bar material constants, cell/DoF counts and partitions. \\
\hyperref[mesh:feelpp:fuel-pellet]{Fuel-pellet fragment} & Height 7 mm; radius 4 mm; fragment angle/shape to be confirmed. Obtain fragment geometry, source amplitude, mechanical masks and thermal markers before generating meshes. \\
\end{FicheMetadataTable}
\paragraph{Next contribution.} Obtain bar constants and pellet geometry/source/constraints, then supply heat/solid CFG/JSON and coupling settings; select Feel++ validation tolerances.

